\section{The least nonnegative Frobenius weight}
\label{sec:detector}
The geometric part of the Lean argument attaches nonnegative local amounts to primes. For a nonabelian extension, one can ask how few Frobenius classes a nonnegative replacement must see while giving every nonlinear character a specified coefficient. There is a sharp pointwise answer.

Let $G$ be a finite group, $H=[G,G]$, $e$ its identity, and $d_\pi=\chi_\pi(e)$ for an irreducible complex character $\chi_\pi$. Expand a real class function satisfying $w(g^{-1})=w(g)$ as $w=\sum_\pi a_\pi\chi_\pi$. Character orthogonality gives real coefficients; conjugate characters have equal coefficients. Write $w_*(g)=\mathbf{1}_{H\setminus\{e\}}(g)$.

\paragraph{Pointwise minimum and an exact stability identity.}
Suppose $w\ge0$ and
\[
             -a_\pi\ \ge\ d_\pi/|G| \qquad(d_\pi>1).
\]
Then $w\ge w_*$ pointwise. Equality is attainable simultaneously at every element. More precisely, put $b_\pi=-a_\pi-d_\pi/|G|$ for $d_\pi>1$. For $h\in H$,
\begin{align*}
 w(h)-w(e)-w_*(h)&=
  \sum_{d_\pi>1} b_\pi\bigl(d_\pi-\operatorname{Re}\chi_\pi(h)\bigr),\\
 \frac1{|H|}\sum_{h\in H}\bigl(w(h)-w(e)-w_*(h)\bigr)
 &=\sum_{d_\pi>1} d_\pi b_\pi .
\end{align*}
Thus the average excess on the commutator subgroup measures precisely how much extra weight went into nonlinear coefficients. If all nonlinear coefficients are strictly negative with arbitrary magnitudes, rescaling reduces to the displayed normalization in the nonabelian case. Every such nonnegative detector must count each nonidentity commutator-subgroup element. The minimization concerns coefficients of one prescribed sign; signed prime weights and tests for selected constituents present different problems.

\paragraph{Proof.}
Linear characters equal $1$ on $H$. In $w(h)-w(e)$ they disappear. A unitary matrix of size $d$ has real trace at most $d$. The regular-character identity gives, for $h\in H$,
\[
 \frac1{|G|}\sum_{d_\pi>1}d_\pi(d_\pi-\operatorname{Re}\chi_\pi(h))
    =1-\mathbf{1}_{\{e\}}(h).
\]
These facts give the first equality and the lower bound. Averaging the representation over $H$ projects onto its $H$-fixed vectors. This subspace is $G$-invariant, and a nonlinear irreducible cannot factor through the abelian $G/H$. Thus its averaged trace is zero, giving the second equality. Finally $w_*=\mathbf{1}_H-\mathbf{1}_{\{e\}}$. The first summand has only linear constituents and the second is the regular character divided by $|G|$. Hence $a_\pi(w_*)=-d_\pi/|G|$ for every nonlinear $\pi$. This proves attainment. In particular, every objective $\sum_g c(g)w(g)$ with all $c(g)>0$ has the unique minimizer $w_*$. The result is a finite-group theorem; analytic use of an arbitrary character expansion also requires control of the relevant Artin factors.

\paragraph{Relation to earlier work.}
The \href{https://github.com/VladimirReshetnikov/ProveIt/blob/5cc267af75a635199ac5d93a7d49cbce521daefc/docs/incoming/grh_deficit_transfer_2026-11-25.zip}{incoming deficit report}, section on optimal Frobenius deficits, already constructs $w_*$ and proves its one-sided zero test. Its minimization fixes a combination of permutation characters and varies the abelian correction. Here every nonlinear coefficient can vary independently subject to its lower bound in absolute value. The \href{https://github.com/VladimirReshetnikov/ProveIt/blob/5cc267af75a635199ac5d93a7d49cbce521daefc/docs/incoming/grh_sparse_singular_moduli.zip}{later Frobenius-weight analysis}, in \texttt{frobenius.tex}, already obtains the strict support obstruction and a sharp normalized minimax bound from the same character identities for arbitrary normal subgroups. The display above records the pointwise excess and its average in the normalization used here. We use this as the baseline for the mixed-sign tests. The \href{https://github.com/VladimirReshetnikov/ProveIt/blob/5cc267af75a635199ac5d93a7d49cbce521daefc/docs/incoming/grh-progression-densities-and-cubic-prime-race.zip}{incoming cubic prime-race report} already proves excursions around the $-2/3$ square-term drift for the trace of $X^3-X-1$ on power-spaced grids. Those oscillations are part of the prior work.

\subsection*{An explicit cubic target}
Let $P$ be a monic irreducible cubic in $\mathbb Z[X]$ with nonsquare discriminant, $E$ its splitting field, and $\tau$ the standard representation of $G=S_3$. Exclude the primes dividing the polynomial discriminant. The support of $w_*$ consists of the two three-cycles, so the minimum weight simply asks whether $P\bmod p$ is irreducible. In the order identity, transposition, three-cycle, $\chi_\tau=(2,0,-1)$ and $\mathrm{sgn}=(1,-1,1)$. Hence
\[
                    3w_* =1+\mathrm{sgn}-\chi_\tau.
\]
Define the first Riesz mean of irreducible primes by
\[
 I_P(x)=3\sum_{\substack{p\le x,\ p\nmid\operatorname{disc}(P)\\
                           P\bmod p\ {\rm irreducible}}}
                  \log p\,(1-p/x).
\]
Under the repository's Dirichlet GRH, the following three conditions are equivalent:
\begin{enumerate}
\item All nontrivial zeros of $L(s,\tau)$ have real part $1/2$.
\item $I_P(x)=x/2+O_P(\sqrt{x})$.
\item For each $\epsilon>0$ and sufficiently large $x$, $I_P(x)\le x/2+C_{P,\epsilon}x^{1/2+\epsilon}$.
\end{enumerate}
This is a single-scale Riesz formulation of the earlier deficit principle. The assumed estimate in the last condition remains open already for $P=X^3-X-1$. The classical degree-two factor for this example comes from $\eta(z)\eta(23z)$; see
\href{https://www.sciencedirect.com/science/article/pii/S0723086923000087}{Khare, La Rosa and Wiese, on splitting fields of $X^n-X-1$}.
For general $S_3$, $\tau$ is induced by a nontrivial character of $A_3$. Hecke theory supplies entireness, the functional equation, and the standard zero count for its $L$-function.

\paragraph{Proof of the criterion.}
Write $U(x)$ for the prime-power version with summands $3w_*(g_p^k)\log p(1-p^k/x)$ and $g_p$ a Frobenius element. Chebyshev's bound $\sum_{p\le y}\log p\ll y$ gives $U-I_P=O(\sqrt{x})$; the exponents at least three contribute $O(x^{1/3}\log x)$. For $\operatorname{Re}s>1$,
\[
 \int_1^\infty U(x)x^{-s-1}\,dx=\frac{A(s)}{s(s+1)},\qquad
 A(s)=-\frac{\zeta'}{\zeta}(s)-\frac{L'}L(s,\mathrm{sgn})
                      +\frac{L'}L(s,\tau)+B(s).
\]
Here $B$ removes the finitely many excluded Euler factors and is holomorphic in $\operatorname{Re}s>0$. At $s=1$ the residue of the transform is $1/2$. After centering by $x/2$, at a real zero $\beta>1/2$ of $L(s,\tau)$ the residue is $m_\beta/[\beta(\beta+1)]>0$. Dirichlet GRH excludes poles from the first two $L$-functions in this region, apart from the accounted-for pole at one.

Suppose (3) holds, fix $a=1/2+\epsilon<1$, and make $V(x)=Cx^a-(U(x)-x/2)$ nonnegative by increasing $C$. If the convergence abscissa $\alpha$ of its Mellin transform were greater than $a$, the transform would have a singularity at the real point $\alpha$ by Landau's theorem. Indeed, Taylor-expanding at a real point just to the right of $\alpha$ and using monotone convergence of the integrals of the powers of $\log x$ rules out regularity there. The meromorphic formula gives either a regular point or a simple pole of strictly negative residue. The latter would send the transform to $-\infty$ as $s\downarrow\alpha$, impossible for $V\ge0$. Thus the transform is holomorphic in $\operatorname{Re}s>a$. Every zero there, real or nonreal, would give a pole, since the kernel $1/[s(s+1)]$ never vanishes. Letting $\epsilon$ decrease to zero and applying the functional equation proves (1).

Under (1), all three factors satisfy their respective critical-line theorems. The Riesz explicit formula gives $U(x)=x/2+O_P(\sqrt{x})$. For details, begin with Mellin inversion on $\operatorname{Re}s=2$ and use the full logarithmic derivatives, adding the omitted local terms back at a cost $O_P(\log x)$. Shift to $\operatorname{Re}s=-1/2$ along heights $T$ at distance $\gg_P1/\log T$ from zero ordinates, possible by the unit-interval zero count $O_P(\log T)$. The log derivatives on the horizontal segments are $O_P(\log^2 T)$ by their local partial-fraction expansion. The functional equations and Stirling's formula give $O_P(\log(2+|t|))$ on the new vertical line away from bounded heights. The crossed nontrivial zero residues have absolute sum $O_P(\sqrt{x})$, since the sum of multiplicities times $(1+|\operatorname{Im}\rho|)^{-2}$ converges. At zero the possible double pole costs $O_P(1+\log x)$. This proves (2), which immediately gives (3). A source for the fixed-degree Artin explicit formula and the discriminant-$-23$ example is
\href{https://www.cs.uleth.ca/~nathanng/RESEARCH/phd.thesis.pdf}{N.~Ng, \emph{Limiting distributions and zeros of Artin $L$-functions}, chapters 3--4}.

For a proposed geometric proof this leaves a precise target: control the total excess of the irreducible-prime mean. The support theorem shows what can be gained by changing the character coefficients within the nonnegative construction. The estimate itself still requires cancellation at arbitrarily large scales.
